\honest{Preface}{Eliezer Yudkowsky}{2015}

This book collects two years of my daily blog posts, and I begin by listing what I did wrong in them. If I saw no mistakes, I would have learned nothing since 2009.

My largest mistake was not realizing that the hard part of learning this way of thinking is practice, not theory. My second was writing about big, distant, abstract problems instead of everyday life, where practice is easier. (The Center for Applied Rationality is working on ``repairing this huge mistake of mine.'') My third was writing too much about belief and too little about action.

My fourth mistake was poor organization, and that one ``at least is correctable'': Rob Bensinger has reordered the posts, without ``trying to rewrite all the actual material.'' \nb{So the book corrects the fourth of the five mistakes. The three largest concern what the posts teach and how, and the preface does not say that any of them was addressed.}

My fifth mistake was open contempt for ideas I thought stupid. I did try to avoid ``Bulverism'': I always discussed an idea before calling it stupid. \nb{Bulverism is C. S. Lewis's name for explaining why someone holds a view before showing that the view is wrong.} My reason: many people take courteous treatment of an idea to mean they will lose no status by saying they believe it, homeopathy for example, and ``derisive laughter by comedians can help people wake up from the dream.'' I still think so, and I give no evidence for it. Today I would write more courteously, ``I think,'' because I now take more seriously the risk of building a community whose expected reaction to low-status outsider views is open mockery. My readers, I am happy to say, have been ``amazingly good'' about not using my rhetoric to bully others.

What was I trying to do? To teach ``a certain valuable way of thinking'' that is ``not taught systematically at all,'' and that some people absorb from books like \textsc{Surely You're Joking, Mr. Feynman}. \nb{The preface gives no support for the claim that it is not taught.} It has to do with science and the experimental method, and with being able to say ``Oops'' and give up on what is not working, in daily life as in the laboratory. Somewhere in your life you are like the people who keep ``shooting yourself in the foot.''

Despite my mistakes, the posts ``appeared to help a surprising number of people a surprising amount.'' I give no number and no example. My explanation is that so little of these skills, and of the mathematics and science behind them, is taught that working through many problems in philosophy and science from many angles can give a glimpse of the ``central rhythm.'' I admit that I ``picked all the wrong examples.'' Twice I answer that ``it is all in the end one thing'': the laws governing distant problems and daily life are the same. I do not argue for this. Five years later, ``it still appears to me that this is better than nothing.''
